% Part: cheat-sheets
% Chapter: quantified-modal-logics
% Section: free-qml
\documentclass[../../../../include/open-logic-section]{subfiles}

\begin{document}

\olfileid{chs}{qml}{fre}

\olsection{Free Quantified Modal Logics}

The systems of Free Quantified Modal Logic, $\Log{FQML}$,
combines free quantification, i.e. positive or negative 
free logic with modal logic. We distinguish 
positive free quantified modal logic, $\Log{PFQML}$,
and negative free quantified modal logic, $\Log{NFQML}$.

Their semantics uses \emph{variable-domain models}.
The differences between negative vs positive semantics
are that:
\begin{itemize}
\item negative semantics constrains the interpretation 
 of predicates to ensure that at a given world, a predicate 
 only applies to objects existing at that world,
 \item the clause for $\eq$ in negative semantics is restricted
 to ensure that at a given true, identity claims are true 
 only of objects that exist at that world. That is, $\eq[\Obj{a}][t]$
 is true at $w$ only if $\Obj{a}$ denotes an object in the domain of $w$. 
\end{itemize}

Their natural deduction system has rules of
\emph{positive free logic} or \emph{negative free logic}
plus the modal rules.

\subsection{Variable-domain models}

In variable-domain models, we have a single \emph{super-domain}
$\mathcal{S}$, but \emph{each world $w$ has its own domain}
$\mDomain{M}(w)$ of objects taken from the super-domain. The world's
domain represents the objects existing at that world.


\begin{defn}[Variable-domain models]
    A \emph{variable-domain model} $\mModel M$ for the language $\Lang{L}$ 
    of quantified modal logic consists of the following elements:
    \begin{enumerate}
    \item \emph{Worlds:} a non-empty set~$\mWorlds{M}$
    \item \emph{Accessibility relation:} a binary relation~$\mAccess{M}$ on~$\mWorlds{M}$
    \item \emph{Super-Domain:} a non-empty set, $\mathcal{S}$
    \item \emph{Domains Function}: a function $\mDomain{M}$ that assigns
    to each world $w$ in $\mWorlds{M}$ a subset $\mDomain{M}(w)$ of $\mathcal{S}$, the \emph{domain of world $w$}
\item \emph{Interpretation of !!{constant}s:} for each !!{constant}~$c$ of
  $\Lang{L}$, !!a{element} of the super-domain, $\mAssign{c}{M} \in \mathcal{S}$
\item \emph{Interpretation of !!{predicate}s:} for each $n$-place
  !!{predicate}~$R$ of $\Lang{L}$ (other than $\eq$), a function 
  from worlds to $n$-place
  relations on the super-domain: $\mAssign{R}{M} \colon W \to \mathcal{S}^n$
\item \emph{Interpretation of !!{function}s:} for each $n$-place
  !!{function}~$f$ of $\Lang{L}$, an $n$-place function from worlds 
  to $n$-place functions on the super-domain: $\mAssign{f}{M} \colon W \to 
  (\mathcal{S}^n \to \mathcal{S})$

\end{enumerate}

We write $\mAssign{R}{M}[w]$ for the value of $\mAssign{R}{M}$ at $w$ (a $n$-place relation) and $\mAssign{f}{M}[w]$ for the value of the 
$\mAssign{f}{M}[w]$ at $w$ (a $n$-place function). 
\end{defn}

\begin{defn}[Negative variable-domain models]
  A variable-domain model is \emph{negative} if it additionally 
  obeys the following constraints:
  \begin{enumerate}
    \item \emph{Interpretation of !!{predicate}s:} the interpretation
    of a $n$-place !!{predicate}~$R$ associates each world $w$ with a
    relation over \emph{the domain of $w$}: $\mAssign{R}{M}[w]\in
    \mDomain{M}(w)$. 
  \item \emph{Interpretation of !!{function}s:} the interpretation 
    an $n$-place !!{function}~$f$ associates each world $w$ with a 
    function \emph{that maps any sequence of objects
    containing some object not in $w$ to an object outside the domain 
    of $w$}. That is, $\mAssign{f}{M}(w)(m_1,\ldots,m_n)\in \mDomain{M}(w)$
    only if $m_1,\ldots,m_n\in \mDomain{M}(w)$.
  \end{enumerate}
\end{defn}

\begin{rem}[Model types]
  As in modal logic, a model $\mModel{M}$ is called a $\Ax{D}$-model, 
  $\Ax{T}$-model, etc., if its accessibility relation is serial,
  reflexive, etc., as in \olref[chs][qml][cla]{tab:five-types}. 
\end{rem}

\subsection{satisfaction}

!!^{variable} assignments are defined as in classical quantified modal 
logic except that we use objects from the super-domain:

\begin{defn}[Variable Assignment]
A \emph{variable assignment}~$s$ for a variable-domain model~$\mModel{M}$ 
is a function which maps each !!{variable} to !!a{element} of
the super-domain~$\mathcal{S}$,
i.e., $s\colon \Var \to \mathcal{S}$.
\end{defn}

The definitions of terms !!{value}s and of 
assignment variants are as in classical quantified modal logic.


Satisfaction is defined as for fixed-domain models, with two differences:
\begin{enumerate}
\item For quantifiers: make their range world-relative,
\item For identity, in the negative semantics: restrict its application
to each world's domain.
\end{enumerate}
Let us look at these differences before stating the full 
definitions of satisfactions. 

\paragraph{Quantifiers} With variable-domain models we restrict the range of 
quantifiers at each world to that world's domain. That includes
the existence predicate:
\begin{itemize}
\item $\qmSat{M}{\Atom{\lfrexists}{t}}[w][s]$ iff 
  $\mValue{t}{M}{w}[s]\in \mDomain{M}(w)$.

\tagitem{prvAll}{%
  $\qmSat{M}{\lforall[x][!B]}[w][s]$ iff for
    every !!{element}~$m \in \mDomain{M}[w]$,
    $\qmSat{M}{!B}[w][\Subst{s}{m}{x}]$.}{}

\tagitem{prvEx}{%
  $\qmSat{M}{\lexists[x][!B]}[w][s]$ iff for at
  least one !!{element}~$m \in \mDomain{M}[w]$,
  $\qmSat{M}{!B}[w][\Subst{s}{m}{x}]$.}{}
\end{itemize}
Note that for a quantified sentence to hold at a world 
$w$ relative to an assignment $s$, we only variants
of $s$ involving objects \emph{of that world's domain}, 
$\mDomain{M}[w]$---not objects of the model's super-domain
 $\mathcal{S}$.

\paragraph{Identity in Negative Free semantics}. For the \emph{positive}
variant of quantified modal logic, the clause for identity
is the same. For the \emph{negative} variant, we restrict it 
as follows:

\begin{itemize}

\item classical, positive free: $\qmSat{M}{\eq[t_1][t_2]}[w][s]$ iff
  $\mValue{t_1}{M}{w}[s] = \mValue{t_2}{M}{w}[s]$.

\item negative free: $\qmSat{M}{\eq[t_1][t_2]}[w][s]$ iff
  $\mValue{t_1}{M}{w}[s]\in \mDomain{M}[w]$ and 
  $\mValue{t_1}{M}{w}[s] = \mValue{t_2}{M}{w}[s]$.

\end{itemize}

Satisfaction has a different identity clause for the 
positive and for the negative variants of free quantified 
modal logic. When we need to disambiguate we write
 $\models_{\Log{PFQML}}$ for the positive semantics and
 $\models_{\Log{NFQML}}$. 

\begin{defn}[Satisfaction in variable-domain models]
\ollabel{defn:satisfaction}
Satisfaction of !!a{formula}~$!A$ a fixed-domain model~$\mModel{M}$
at a world $w$ in $\mWorlds{M}$ relative to !!a{variable} assignment~$s$, in symbols:
$\qmSat{M}{!A}[w][s]$, is defined recursively as follows. We write
$\qmSat/{M}{!A}[w][s]$ to mean ``not $\qmSat{M}{!A}[w][s]$.'' 
When we need to specify the positive or negative variant of 
the definition, we use  When we need to disambiguate we write
 $\models_{\Log{PFQML}}$ for the positive semantics and
 $\models_{\Log{NFQML}}$. 
\begin{enumerate}
\tagitem{prvFalse}{%
  \indcase{!A}{\lfalse}{$\qmSat/{M}{\indfrm}[w][s]$.}}{}

\tagitem{prvTrue}{%
  \indcase{!A}{\ltrue}{$\qmSat{M}{\indfrm}[w][s]$.}}{}

\item \indcase{!A}{\Atom{R}{t_1, \dots, t_n}}{$\qmSat{M}{\indfrm}[w][s]$
  iff $\langle \mValue{t_1}{M}{w}[s], \dots, \mValue{t_n}{M}{w}[s] \rangle \in
  \mAssign{R}{M}[w]$.}

\item \indcase{!A}{\Atom{\lfrexists}{t}}{$\qmSat{M}{\indfrm}[w][s]$ iff 
  $\mValue{t}{M}{w}[s]\in \mDomain{M}(w)$.}

\item \indcase{!A}{\eq[t_1][t_2]}{$\qmSat{M}{\indfrm}[w][s]$ iff
  \begin{itemize}
    \item positive variant: $\mValue{t_1}{M}{w}[s] = \mValue{t_2}{M}{w}[s]$,
    \item negative variant:  $\mValue{t_1}{M}{w}[s]\in \mDomain{M}[w]$ and 
  $\mValue{t_1}{M}{w}[s] = \mValue{t_2}{M}{w}[s]$.
  \end{itemize}
 }

\tagitem{prvNot}{%
  \indcase{!A}{\lnot !B}{$\qmSat{M}{\indfrm}[w][s]$ iff
    $\qmSat/{M}{!B}[w][s]$.}}{}

\tagitem{prvAnd}{%
  \indcase{!A}{(!B \land !C)}{$\qmSat{M}{\indfrm}[w][s]$ iff 
  $\qmSat{M}{!B}[w][s]$ and $\qmSat{M}{!C}[w][s]$.}}{}

\tagitem{prvOr}{%
  \indcase{!A}{(!B \lor !C)}{$\qmSat{M}{\indfrm}[w][s]$ iff
    $\qmSat{M}{!B}[w][s]$ or $\qmSat{M}{!C}[w][s]$ (or both).}}{}

\tagitem{prvIf}{%
  \indcase{!A}{(!B \lif !C)}{$\qmSat{M}{\indfrm}[w][s]$ iff 
    $\qmSat/{M}{!B}[w][s]$ or $\qmSat{M}{!C}[w][s]$ (or both).}}{}

\tagitem{prvIff}{%
  \indcase{!A}{(!B \liff !C)}{$\qmSat{M}{\indfrm}[w][s]$ iff either
    both $\qmSat{M}{!B}[w][s]$ and $\qmSat{M}{!C}[w][s]$, or neither
    $\qmSat{M}{!B}[w][s]$ nor $\qmSat{M}{!C}[w][s]$.}}{}

\tagitem{prvAll}{%
  \indcase{!A}{\lforall[x][!B]}{$\qmSat{M}{\indfrm}[w][s]$ iff for
    every !!{element}~$m \in \mDomain{M}(w)$,
    $\qmSat{M}{!B}[w][\Subst{s}{m}{x}]$.}}{}

\tagitem{prvEx}{%
  \indcase{!A}{\lexists[x][!B]}{$\qmSat{M}{\indfrm}[w][s]$ iff for at
  least one !!{element}~$m \in \mDomain{M}(w)$,
  $\qmSat{M}{!B}[w][\Subst{s}{m}{x}]$.}}{}

\tagitem{prvBox}{%
  \indcase{!A}{\Box !B}{$\qmSat{M}{\indfrm}[w][s]$ iff
    $\qmSat{M}{!B}[w'][s]$ for all $w' \in \mWorlds{M}$ with
    $\mAccess{M}ww'$.}}{}

\tagitem{prvDiamond}{%
    \indcase{!A}{\Diamond !B}{$\qmSat{M}{\indfrm}[w][s]$ iff
      $\qmSat{M}{!B}[w'][s]$ for at least one $w' \in \mWorlds{M}$
      with $\mAccess{M}ww'$.}}{}

\end{enumerate}
\end{defn}

Sentences and truth at a world are defined as in classical quantified
modal logic.

\subsection{Natural Deduction for \Log{FQML}}

We use two systems for free quantified modal logic:
either positive \Log{PFQML} or negative \Log{NFQML}. 

You are allowed to use the following rules, detailed
in the relevant sections of these cheatsheets:

\begin{enumerate}
  \item Propositional Logic. All rules of propositional logic.
  \item Free Quantification. All rules of quantifiers in free
  logics (they are the same inpositive and negative).
  \item The positive or negative free logic rules
  for identity and atomic predication, respectively. 
  \begin{enumerate} 
    \item Positive (\Log{PFQML}): \Intro{\eq} and \Elim{\eq} as in
    first order logic. No constraint on atomic predication.
    \item Negative (\Log{PFQML}): the weaker \Intro{\eq}$_{\Log{FL}}$ rule, and
    \Elim{\eq}. Constraint on atomic predication: \Log{NFL} 
    rules that allow us to infer $\lfrexists t$ from 
    $\Atom{R}{\ldots,t,\ldots}$ and $\eq[t][t']$.
  \end{enumerate}
  \item Existence. The definition of the existence predicate
  $\lfrexists$ (below). These are derived rules in free logic systems.
  \item Modality. All rules of modal logics. Depending on how the
    accessibility 
    relation of our models are constrained:
    \begin{enumerate}
      \item By default, \Log{PFQML} and \Log{PFQML} use modal rules
      from the basic modal system \Log{K}.
      \item Stronger systems are \Log{FQML+X}, for instance
      \Log{NFQML+T} is negative free quantified modal logic with modal rules
      from system \Log{T}. Its models are variable-domain models where
      the accessibility relation is reflexive.
      \item Sider's "Variable-Domain Quantified Modal Logic" \Log{SQML}
      is \Log{PFQML+S5}, positive free quantified modal logic with 
      modal rules for system \Log{S5}. Its models are variable-domain models where
      the accessibility relation is an equivalence relation (reflexive, 
      symmetric and transitive), and it uses the positive 
      variant of satisfaction clauses.
    \end{enumerate}
\end{enumerate}

\begin{defish}
    \emph{Definition of $\lfrexists$}

    \begin{center}        
    \smallskip
    \begin{tabular}{cc}
    \AxiomC{$\lexists[v][\eq[c][v]]$}
    \RightLabel{Def~$\lfrexists$}
    \UnaryInfC{$\lfrexists c$}
    \DisplayProof
    &
    \AxiomC{$\lfrexists c$}
    \RightLabel{Def~$\lfrexists$}
    \UnaryInfC{$\lexists[v][\eq[c][v]]$}
    \DisplayProof
    \end{tabular}
    \end{center}

\end{defish}

\end{document}